\documentclass[a4paper, serif, 10pt]{moderncv}

%% moderncv
% color and style
\moderncvcolor{black}
\moderncvstyle{banking}

%% page layout/margins
\usepackage[top=2.5cm, bottom=2.5cm, left=1.5cm, right=1.5cm]{geometry}

\nopagenumbers{}

%% Babel config for Indonesian (Bahasa Indonesia) language
% ref:
% - https://latex3.github.io/babel/guides/locale-indonesian.html
\usepackage[indonesian]{babel}

%% fontspec
\usepackage{fontspec}

\IfFontExistsTF{Linux Libertine}{
  \setmainfont[Ligatures={TeX, Common}, Numbers=OldStyle]{Linux Libertine}
}{}
\IfFontExistsTF{Linux Libertine O}{
  \setmainfont[Ligatures={TeX, Common}, Numbers=OldStyle]{Linux Libertine O}
}{}

%% CTeX
% CJK support, used to show CJK characters in the resume
%
% - heading=false: disable ctex heading
% - punct=kaiming: use kaiming punctuations styles for CJK
% - scheme=plain: use plain scheme, do not override \normalsize font size
% - space=auto: space settings for CJK characters
%
% CJK fonts are resolved per locale language, e.g. Noto Serif/Sans CJK JP
% for Japanese resumes, with fallback to the next region if not installed
%
% ref:
% - http://ctan.mirrorcatalogs.com/language/chinese/ctex/ctex.pdf
\usepackage[UTF8, heading=false, punct=kaiming, scheme=plain, space=auto]{ctex}

\IfFontExistsTF{Noto Serif CJK SC}{
  \setCJKmainfont{Noto Serif CJK SC}
}{}
\IfFontExistsTF{Noto Sans CJK SC}{
  \setCJKsansfont{Noto Sans CJK SC}
}{}

%% line spacing
\usepackage{setspace}
\setstretch{1.125}

%% URL styling - use normal text instead of monospace
\urlstyle{same}

%% Patch \httplink and \httpslink to handle full URLs
% The original moderncv macros always prepend http:// or https:// to URLs.
% This causes issues when the URL already has a protocol (e.g., https://example.com)
% resulting in malformed URLs like https://https://example.com.
% This patch checks if the URL already contains "://" and uses it directly if so.
% Note: We use \str_set:Nx (x-expansion) to expand macros like \@homepage first.
\ExplSyntaxOn
\str_new:N \g_yamlresume_url_str

\RenewDocumentCommand{\httpslink}{O{}m}{
  \str_set:Nx \g_yamlresume_url_str {#2}
  \str_if_in:NnTF \g_yamlresume_url_str {://}
    { \url{#2} }
    { \url{https://#2} }
}

\RenewDocumentCommand{\httplink}{O{}m}{
  \str_set:Nx \g_yamlresume_url_str {#2}
  \str_if_in:NnTF \g_yamlresume_url_str {://}
    { \url{#2} }
    { \url{http://#2} }
}
\ExplSyntaxOff

\name{Raka Aditya Pratama}{}
\title{Cybersecurity Engineer | Spesialis Keamanan Jaringan dan Cloud}
\phone[mobile]{+62 812-3456-7890}
\email{raka.pratama@cybermail.id}
\homepage{https://www.rakaaditya.id}

\address{Jl. Cikini Raya No. 45, Jakarta Pusat, DKI Jakarta, Indonesia, 10330}{}{}

\extrainfo{{\small \faLinkedin}\ \href{https://www.linkedin.com/in/rakaadityapratama}{@rakaadityapratama} {} {} {} • {} {} {} 
{\small \faGithub}\ \href{https://github.com/rakaaditya}{@rakaaditya}}

\begin{document}

\maketitle

\section{Informasi Dasar}

\cvline{}{Cybersecurity Engineer dengan pengalaman lebih dari 6 tahun dalam merancang, mengoperasikan, dan meningkatkan kontrol keamanan untuk lingkungan enterprise, cloud, dan aplikasi web. Terampil dalam threat hunting, incident response, pengujian penetrasi, serta hardening infrastruktur dan identitas.
%
Fokus pada pengurangan risiko siber secara terukur melalui otomasi, kolaborasi lintas tim, dan kepatuhan terhadap standar ISO 27001, NIST CSF, dan PCI DSS. Berpengalaman memimpin investigasi insiden, membangun runbook respons, dan melatih tim internal dalam praktik keamanan yang aman.}
\section{Pendidikan}

\cventry{Agu 2013–Jul 2017}
        {Sarjana, Ilmu Komputer, Nilai: 3.72}
        {Universitas Indonesia}
        {\url{https://www.ui.ac.id}}
        {}
        {\begin{itemize}
\item Menguasai dasar-dasar keamanan informasi, kriptografi, dan forensik digital melalui perkuliahan dan proyek laboratorium.
\item Mengembangkan aplikasi berbasis web dan melakukan analisis kerentanan dasar sebagai bagian dari tugas akhir.
\item Aktif dalam kelompok studi keamanan siber kampus dan membantu penyelenggaraan workshop Capture The Flag.
\end{itemize}
\textbf{Mata Kuliah}: Keamanan Informasi, Jaringan Komputer, Kriptografi, Sistem Operasi, Basis Data, Pemrograman Lanjut, Forensik Digital, Arsitektur Perangkat Lunak}
\section{Pengalaman Kerja}

\cventry{Mar 2022–Sekarang}
        {Senior Cybersecurity Engineer}
        {PT Nusantara Cyber Solutions}
        {\url{https://www.nusantaracyber.co.id}}
        {}
        {\begin{itemize}
\item Memimpin tim blue team dalam pemantauan keamanan 24/7, threat hunting, dan investigasi insiden untuk klien enterprise di sektor keuangan dan energi.
\item Merancang dan mengimplementasikan otomasi incident response berbasis SOAR yang memangkas waktu triase alert hingga 40\%.
\item Melakukan pengujian penetrasi aplikasi web, API, dan infrastruktur cloud serta menyusun rekomendasi remediasi yang dapat ditindaklanjuti.
\item Mengembangkan runbook, playbook, dan matriks eskalasi insiden yang selaras dengan NIST CSF dan ISO 27001.
\item Melatih analis junior dan tim TI klien mengenai praktik secure configuration, phishing awareness, dan manajemen kerentanan.
\end{itemize}
\textbf{Kata Kunci}: Incident Response, Threat Hunting, SIEM, ISO 27001}

\cventry{Jun 2019–Feb 2022}
        {Cybersecurity Engineer}
        {PT Bank Digital Sentosa}
        {\url{https://www.bankdigitalsentosa.co.id}}
        {}
        {\begin{itemize}
\item Mengelola dan menyempurnakan aturan deteksi pada SIEM, mendukung monitoring transaksi digital dan aktivitas pengguna berisiko.
\item Melakukan vulnerability assessment dan penetration testing pada aplikasi mobile banking, portal web, dan layanan API.
\item Berkontribusi dalam program kepatuhan PCI DSS dan ISO 27001, termasuk pengumpulan bukti audit dan remediasi temuan.
\item Mengimplementasikan network segmentation, hardening server, dan kontrol akses berbasis prinsip least privilege.
\item Bekerja sama dengan tim fraud dan risiko untuk mendeteksi serta merespons indikasi penipuan siber secara cepat.
\end{itemize}
\textbf{Kata Kunci}: PCI DSS, Vulnerability Assessment, Network Security, Mobile Security}

\cventry{Agu 2017–Mei 2019}
        {Security Operations Analyst}
        {PT Teknologi Aman Indonesia}
        {\url{https://www.teknologiaman.id}}
        {}
        {\begin{itemize}
\item Memantau alert keamanan dari SIEM, firewall, IDS/IPS, dan endpoint protection untuk mendeteksi aktivitas mencurigakan.
\item Melakukan triase, klasifikasi, dan eskalasi insiden sesuai SLA serta mendokumentasikan temuan dalam tiket investigasi.
\item Membantu analisis malware dasar, pengumpulan bukti digital, dan penyusunan laporan insiden mingguan.
\item Mengotomatisasi tugas pengumpulan log dan enrichment IOC menggunakan skrip Python dan API keamanan.
\item Mendukung penyusunan prosedur operasional keamanan dan knowledge base untuk tim SOC.
\end{itemize}
\textbf{Kata Kunci}: SOC, IDS/IPS, Malware Analysis, Python}
\section{Bahasa}

\cvline{Indonesia}{Bahasa Ibu / Bilingual \hfill \textbf{Kata Kunci}: Bahasa ibu}
\cvline{Inggris}{Tingkat Profesional Penuh \hfill \textbf{Kata Kunci}: TOEFL iBT 105, IELTS 7.5}
\section{Keahlian}

\cvline{Keamanan Jaringan}{Ahli \hfill \textbf{Kata Kunci}: Firewall, IDS/IPS, VPN, Network Segmentation, Zero Trust}
\cvline{Pengujian Penetrasi}{Mahir \hfill \textbf{Kata Kunci}: Burp Suite, Metasploit, Nmap, OWASP Top 10, Kali Linux}
\cvline{Keamanan Cloud}{Mahir \hfill \textbf{Kata Kunci}: AWS Security Hub, IAM, Azure Defender, CSPM, Kubernetes Security}
\cvline{Incident Response}{Ahli \hfill \textbf{Kata Kunci}: SIEM, SOAR, Threat Hunting, Digital Forensics, Malware Analysis}
\cvline{Tata Kelola dan Kepatuhan}{Mahir \hfill \textbf{Kata Kunci}: ISO 27001, PCI DSS, NIST CSF, Risk Assessment, Audit}
\section{Penghargaan}

\cventry{Nov 2023}
        {PT Nusantara Cyber Solutions}
        {Penghargaan Inovasi Keamanan Terbaik}
        {}
        {}
        {Diberikan atas perancangan platform otomasi incident response yang mengurangi waktu deteksi dan respons insiden sebesar 40\%.}
\section{Sertifikat}

\cventry{Jun 2023}
        {ISC2}
        {Certified Information Systems Security Professional (CISSP)}
        {\url{https://www.isc2.org/certifications/cissp}}
        {}
        {}

\cventry{Sep 2021}
        {Offensive Security}
        {Offensive Security Certified Professional (OSCP)}
        {\url{https://www.offsec.com/certifications/oscp/}}
        {}
        {}

\cventry{Mar 2022}
        {Amazon Web Services}
        {AWS Certified Security - Specialty}
        {\url{https://aws.amazon.com/certification/certified-security-specialty/}}
        {}
        {}
\section{Publikasi}

\cventry{Mar 2023}
        {Deteksi Anomali Jaringan Menggunakan Pembelajaran Mesin untuk Lingkungan Cloud}
        {Jurnal Keamanan Siber Indonesia}
        {\url{https://jurnal.keamanansiber.id/volume-7/no-1/deteksi-anomali-jaringan}}
        {}
        {\begin{itemize}
\item Membahas pendekatan deteksi anomali berbasis pembelajaran mesin untuk memantau lalu lintas jaringan cloud.
\item Mengevaluasi performa model terhadap dataset lalu lintas normal dan serangan yang disimulasikan.
\item Memberikan rekomendasi implementasi untuk tim keamanan yang mengelola lingkungan hybrid cloud.
\end{itemize}}
\section{Referensi}

\cventry{\emaillink[intan.prameswari@ui.ac.id]{intan.prameswari@ui.ac.id}}
        {Dosen Keamanan Informasi, Universitas Indonesia}
        {Dr. Intan Prameswari}
        {+62 811-9876-5432}
        {}
        {Raka menunjukkan kemampuan analisis keamanan yang kuat, etika kerja yang tinggi, dan kepemimpinan yang baik dalam proyek incident response. Ia mampu menghubungkan temuan teknis dengan kebutuhan risiko bisnis secara jelas.}
\section{Proyek}

\cventry{Jan 2023–Okt 2023}
        {Platform internal untuk mengotomatisasi triase, enrichment, dan eskalasi alert keamanan.}
        {Platform Otomasi Incident Response}
        {\url{https://github.com/rakaaditya/ir-automation}}
        {}
        {\begin{itemize}
\item Mengintegrasikan SIEM, threat intelligence, dan sistem tiket untuk mempercepat alur investigasi insiden.
\item Menyediakan playbook otomatis untuk phishing, malware, dan akses tidak sah.
\item Mengurangi waktu respons awal dari rata-rata 25 menit menjadi 15 menit.
\end{itemize}
\textbf{Kata Kunci}: Otomasi, SOAR, Threat Intelligence, Python}

\cventry{Apr 2021–Des 2021}
        {Alat pemindaian kerentanan aplikasi web untuk mendeteksi konfigurasi tidak aman dan celah umum.}
        {Pemindaian Kerentanan Aplikasi Web}
        {\url{https://github.com/rakaaditya/web-vuln-scanner}}
        {}
        {\begin{itemize}
\item Mengotomatisasi pemeriksaan header keamanan, cookie, dan konfigurasi TLS.
\item Mendeteksi indikasi kerentanan OWASP Top 10 melalui pengujian non-destruktif.
\item Menghasilkan laporan temuan yang dapat diekspor untuk kebutuhan remediasi dan audit.
\end{itemize}
\textbf{Kata Kunci}: OWASP, Web Security, Python, Automation}
\section{Minat}

\cvline{Keamanan Siber}{Capture The Flag, Riset Kerentanan, Threat Intelligence}
\cvline{Kegiatan Luar Ruang}{Mendaki, Bersepeda, Fotografi}
\section{Kegiatan Sukarela}

\cventry{Feb 2020–Sekarang}
        {Mentor Keamanan Aplikasi}
        {Komunitas Keamanan Siber Indonesia}
        {\url{https://www.keamanansiber.id/komunitas}}
        {}
        {\begin{itemize}
\item Membimbing peserta workshop tentang secure coding, pengujian penetrasi web, dan manajemen kerentanan.
\item Membantu menyusun materi pelatihan praktis untuk pemula dan profesional TI.
\item Berperan sebagai juri dan fasilitator dalam kompetisi Capture The Flag tingkat nasional.
\end{itemize}}
    
\end{document}